\documentclass[12pt]{article}
\usepackage[UTF8]{ctex}  % 中文支持

\usepackage[utf8]{inputenc}
\usepackage{amsmath}
\usepackage{amssymb}
\usepackage{geometry}
\usepackage{graphicx}
\usepackage{amsmath, amssymb}
\usepackage{geometry}
\usepackage{tikz}

\geometry{a4paper, margin=1in}

% \title{Homework 1206}
% \author{Liu}
% \date{\today}

\begin{document}

\section*{公因数与公倍数}
\subsection{写出下面各组数的最大公因数}
\begin{enumerate}
    \item 36和24
    \item 39和13
    \item 31和32
    \item 30和26
    \item 9和16
    \item 91和13
\end{enumerate}
\subsection{写出下面各组数的最小公因数}
\begin{enumerate}
    \item 18和24
    \item 9和7
    \item 17和51
    \item 21和14
    \item 9和10
    \item 65和13
\end{enumerate}
\subsection{应用题}
小明家客厅长5.6米，宽4.8米，他准备选用正方形瓷砖铺地面（不切割瓷砖），
至少需要（）块瓷砖，选用的瓷砖变边长是（）米。
\begin{enumerate}
     \item 如果 $A = 2 \times 2 \times 5$, $B = 2 \times 3 \times 5$，那么 $A$ 和 $B$ 的最小公倍数是（\quad）。 \newline \textbf{A.} 10 \quad \textbf{B.} 30 \quad \textbf{C.} 60 \quad \textbf{D.} 120 
     \item $a$ 和 $a - 1$（$a$ 和 $a - 1$ 都是非零自然数）的最小公倍数是（\quad）。\newline \textbf{A.} $a$ \quad \textbf{B.} $a - 1$ \quad \textbf{C.} $a(a - 1)$ \quad \textbf{D.} 无法确定 \item 学校在一面墙的正方形区域展出同学们的艺术节美术作品。作品规格是长 $42\text{cm}$，宽 $30\text{cm}$，作品刚好占满展区（作品间的间隙忽略不计）。这次至少展出（\quad）幅作品。 \end{enumerate}

\end{document}